\documentclass[11pt]{article}
\newcommand{\ListaTitulo}{Lista 11: VA contínua, esperança, variância e FDA}
% Preâmbulo compartilhado: MATF14, Listas de Exercícios 2026
% Usado por ListaNN.tex e GabaritoNN.tex via % Preâmbulo compartilhado: MATF14, Listas de Exercícios 2026
% Usado por ListaNN.tex e GabaritoNN.tex via % Preâmbulo compartilhado: MATF14, Listas de Exercícios 2026
% Usado por ListaNN.tex e GabaritoNN.tex via \input{preamble}

\usepackage[utf8]{inputenc}
\usepackage[T1]{fontenc}
\usepackage[brazil]{babel}
\usepackage{fullpage}
\usepackage{amsmath,amssymb}
\usepackage{enumitem}
\usepackage{xcolor}
\usepackage{fancyhdr}
\usepackage{titlesec}
\usepackage{listings}
\usepackage{hyperref}
\usepackage{booktabs}
\usepackage{graphicx}
\usepackage{tikz}

\definecolor{matf14azul}{HTML}{0D3B54}
\definecolor{matf14dourado}{HTML}{C9971E}

\titleformat{\section}{\normalfont\Large\bfseries\color{matf14azul}}{\thesection}{1em}{}
\titleformat{\subsection}{\normalfont\large\bfseries\color{matf14azul}}{\thesubsection}{1em}{}

\providecommand{\ListaTitulo}{Lista de Exercícios}

\setlength{\headheight}{14pt}
\pagestyle{fancy}
\fancyhf{}
\lhead{\small MATF14: Estatística Econômica I}
\rhead{\small \ListaTitulo}
\cfoot{\thepage}
\renewcommand{\headrulewidth}{0.4pt}

\lstset{
  basicstyle=\ttfamily\small,
  breaklines=true,
  frame=single,
  columns=fullflexible,
  backgroundcolor=\color{gray!8},
  commentstyle=\color{matf14dourado},
  keywordstyle=\color{matf14azul}\bfseries,
  language=R
}

\newcounter{questaocounter}
\newenvironment{questao}[1]{%
  \stepcounter{questaocounter}%
  \bigskip\noindent\textbf{Questão #1.}\ %
}{\par}

\newcommand{\cabecalho}[2]{%
  \begin{center}
    {\Large\bfseries\color{matf14azul} #1}\\[0.3em]
    {\large #2}
  \end{center}
  \vspace{0.5em}
  \noindent\rule{\linewidth}{0.6pt}
  \vspace{0.5em}
}

\newenvironment{solucao}{%
  \par\smallskip\noindent\textit{\color{matf14dourado!80!black}Solução.}\ %
}{\hfill$\blacksquare$\par}


\usepackage[utf8]{inputenc}
\usepackage[T1]{fontenc}
\usepackage[brazil]{babel}
\usepackage{fullpage}
\usepackage{amsmath,amssymb}
\usepackage{enumitem}
\usepackage{xcolor}
\usepackage{fancyhdr}
\usepackage{titlesec}
\usepackage{listings}
\usepackage{hyperref}
\usepackage{booktabs}
\usepackage{graphicx}
\usepackage{tikz}

\definecolor{matf14azul}{HTML}{0D3B54}
\definecolor{matf14dourado}{HTML}{C9971E}

\titleformat{\section}{\normalfont\Large\bfseries\color{matf14azul}}{\thesection}{1em}{}
\titleformat{\subsection}{\normalfont\large\bfseries\color{matf14azul}}{\thesubsection}{1em}{}

\providecommand{\ListaTitulo}{Lista de Exercícios}

\setlength{\headheight}{14pt}
\pagestyle{fancy}
\fancyhf{}
\lhead{\small MATF14: Estatística Econômica I}
\rhead{\small \ListaTitulo}
\cfoot{\thepage}
\renewcommand{\headrulewidth}{0.4pt}

\lstset{
  basicstyle=\ttfamily\small,
  breaklines=true,
  frame=single,
  columns=fullflexible,
  backgroundcolor=\color{gray!8},
  commentstyle=\color{matf14dourado},
  keywordstyle=\color{matf14azul}\bfseries,
  language=R
}

\newcounter{questaocounter}
\newenvironment{questao}[1]{%
  \stepcounter{questaocounter}%
  \bigskip\noindent\textbf{Questão #1.}\ %
}{\par}

\newcommand{\cabecalho}[2]{%
  \begin{center}
    {\Large\bfseries\color{matf14azul} #1}\\[0.3em]
    {\large #2}
  \end{center}
  \vspace{0.5em}
  \noindent\rule{\linewidth}{0.6pt}
  \vspace{0.5em}
}

\newenvironment{solucao}{%
  \par\smallskip\noindent\textit{\color{matf14dourado!80!black}Solução.}\ %
}{\hfill$\blacksquare$\par}


\usepackage[utf8]{inputenc}
\usepackage[T1]{fontenc}
\usepackage[brazil]{babel}
\usepackage{fullpage}
\usepackage{amsmath,amssymb}
\usepackage{enumitem}
\usepackage{xcolor}
\usepackage{fancyhdr}
\usepackage{titlesec}
\usepackage{listings}
\usepackage{hyperref}
\usepackage{booktabs}
\usepackage{graphicx}
\usepackage{tikz}

\definecolor{matf14azul}{HTML}{0D3B54}
\definecolor{matf14dourado}{HTML}{C9971E}

\titleformat{\section}{\normalfont\Large\bfseries\color{matf14azul}}{\thesection}{1em}{}
\titleformat{\subsection}{\normalfont\large\bfseries\color{matf14azul}}{\thesubsection}{1em}{}

\providecommand{\ListaTitulo}{Lista de Exercícios}

\setlength{\headheight}{14pt}
\pagestyle{fancy}
\fancyhf{}
\lhead{\small MATF14: Estatística Econômica I}
\rhead{\small \ListaTitulo}
\cfoot{\thepage}
\renewcommand{\headrulewidth}{0.4pt}

\lstset{
  basicstyle=\ttfamily\small,
  breaklines=true,
  frame=single,
  columns=fullflexible,
  backgroundcolor=\color{gray!8},
  commentstyle=\color{matf14dourado},
  keywordstyle=\color{matf14azul}\bfseries,
  language=R
}

\newcounter{questaocounter}
\newenvironment{questao}[1]{%
  \stepcounter{questaocounter}%
  \bigskip\noindent\textbf{Questão #1.}\ %
}{\par}

\newcommand{\cabecalho}[2]{%
  \begin{center}
    {\Large\bfseries\color{matf14azul} #1}\\[0.3em]
    {\large #2}
  \end{center}
  \vspace{0.5em}
  \noindent\rule{\linewidth}{0.6pt}
  \vspace{0.5em}
}

\newenvironment{solucao}{%
  \par\smallskip\noindent\textit{\color{matf14dourado!80!black}Solução.}\ %
}{\hfill$\blacksquare$\par}


\begin{document}

\cabecalho{Lista de Exercícios 11}{Variável aleatória contínua, densidade, esperança, variância e FDA (Cap. 4, Aulas 24--26)}

\noindent \textit{Justifique todas as respostas. As resoluções são discutidas em sala e nos horários de atendimento; não são publicadas neste material.}

\begin{questao}{1} \textbf{(Densidade válida: tempo de vida de um equipamento)}

O tempo de vida útil $X$ (em anos) de um equipamento industrial tem densidade
$$
f(x) = k\,x, \quad 0\le x\le 4, \qquad f(x)=0 \text{ caso contrário.}
$$

\begin{enumerate}[label=(\alph*)]
  \item Determine o valor de $k$ para que $f$ seja uma densidade de probabilidade válida.
  \item Calcule $P(X\le 2)$.
  \item Calcule $P(1\le X\le 3)$.
\end{enumerate}
\end{questao}

\begin{questao}{2} \textbf{(Esperança e variância: receita de um posto de vendas)}

A receita diária $X$ (em milhares de reais) de um pequeno posto de vendas de artesanato tem
densidade $f(x) = \frac{1}{8}$ para $2\le x\le 10$ (Uniforme), e $f(x)=0$ fora desse intervalo.

\begin{enumerate}[label=(\alph*)]
  \item Calcule $E(X)$ usando a fórmula geral $E(X)=\int x f(x)\,dx$ (não a fórmula pronta da
    Uniforme, o objetivo aqui é praticar a integral).
  \item Calcule $\mathrm{Var}(X)$, também pela definição via integral.
  \item Confira os dois resultados usando as fórmulas prontas $E(X)=\frac{a+b}{2}$ e
    $\mathrm{Var}(X)=\frac{(b-a)^2}{12}$.
\end{enumerate}
\end{questao}

\begin{questao}{3} \textbf{(FDA a partir da densidade: desconto em uma promoção)}

O percentual de desconto $X$ (entre 0 e 20\%) oferecido em um cupom promocional aleatório tem
densidade $f(x) = \frac{3}{8000}x^2$, $0\le x\le 20$.

\begin{enumerate}[label=(\alph*)]
  \item Encontre a FDA $F(x)$ para $0\le x\le 20$ (calcule a primitiva de $f$).
  \item Use $F$ para calcular a probabilidade de um cliente receber um desconto de pelo menos 15\%.
  \item Calcule a mesma probabilidade do item (b) diretamente por integração de $f$, e confirme
    que os dois métodos dão o mesmo resultado.
\end{enumerate}
\end{questao}

\begin{questao}{4} \textbf{(Propriedades: conversão de moeda)}

O preço $X$ (em dólares) de um insumo importado tem $E(X)=50$ e $dp(X)=8$. Um importador converte
para reais a uma cotação fixa de R\$ 5,20/dólar mais uma taxa fixa de desembaraço de R\$ 30:
$C = 5{,}20X + 30$.

\begin{enumerate}[label=(\alph*)]
  \item Calcule $E(C)$.
  \item Calcule $dp(C)$.
  \item Se a cotação subir para R\$ 5,50/dólar (mantendo a mesma taxa fixa), o desvio padrão de
    $C$ aumenta, diminui ou permanece igual? Justifique sem recalcular tudo, só usando a
    propriedade de $\mathrm{Var}(aX+b)$.
\end{enumerate}
\end{questao}

\begin{questao}{5} \textbf{(Dedução: $P(X=x)=0$ para VA contínua)}

Seja $X$ uma VA contínua com FDA $F$ contínua (sem saltos) em todo $\mathbb{R}$.

\begin{enumerate}[label=(\alph*)]
  \item Para $\varepsilon>0$ pequeno, escreva $P(x-\varepsilon < X \le x)$ em termos de $F$.
  \item Argumente que, como $F$ é contínua, $\displaystyle\lim_{\varepsilon\to0^+}
    \big[F(x)-F(x-\varepsilon)\big] = 0$.
  \item Usando (a) e (b), justifique por que $P(X=x)=0$ para todo $x$, quando $X$ é uma VA
    contínua. (Note que isso \emph{não} valeria se $F$ tivesse um salto em $x$, é exatamente essa
    ausência de saltos, no caso contínuo, que garante o resultado.)
\end{enumerate}
\end{questao}

\bigskip
\noindent\textit{A questão a seguir não tem resposta única e não é cobrada em prova no mesmo
formato: é para discussão em grupo ou por escrito, antes da próxima aula.}

\begin{questao}{6 (discussão)} \textbf{(Sem resposta única)}

Se $P(X=x)=0$ para todo valor específico $x$ de uma variável aleatória contínua, escreva, para um
colega que ainda não cursou esta disciplina, uma explicação de por que isso não é uma
contradição com o fato de $X$ certamente assumir \emph{algum} valor quando o experimento
acontece. Evite fórmulas, o objetivo é uma explicação em português comum.
\end{questao}


\bigskip
\section*{Exercícios complementares}

\begin{enumerate}[label=\textbf{C\arabic*.}, leftmargin=*]
\item (Aquecimento de cálculo) Calcule as derivadas: $8x^{11}$; $5 + x + 3x^2$; $e^{-3x}$; $e^{7x^2 - 2x}$. Calcule as integrais: $\int_0^1 7\,dx$; $\int_1^2 x^3\,dx$; $\int_0^2 (x^2 - 3x + 5)\,dx$; $\int_0^3 e^{2x}\,dx$.
\item (Triangular em $[0, 1]$) Seja $f(x) = Cx$ para $0 \leq x \leq 1/2$, $f(x) = C(1 - x)$ para $1/2 < x \leq 1$ e $f(x) = 0$ fora.
\begin{enumerate}[label=(\alph*)]
  \item Encontre $C$.
  \item Esboce o gráfico de $f$.
  \item Calcule $P(X \leq 1/2)$ e $P(1/4 \leq X \leq 3/4)$.
\end{enumerate}
\item Seja $f(x) = C(1 - x^2)$ para $-1 \leq x \leq 1$ e zero fora.
\begin{enumerate}[label=(\alph*)]
  \item Encontre $C$.
  \item Calcule $E(X)$.
  \item Calcule $Var(X)$.
\end{enumerate}
\item (Exponencial) Para $\lambda > 0$, seja $f(x) = \lambda e^{-\lambda x}$ para $x \geq 0$ e zero para $x < 0$.
\begin{enumerate}[label=(\alph*)]
  \item Encontre a função de distribuição acumulada de $X$.
  \item Calcule $P(X > 10)$.
  \item Encontre a mediana de $X$.
  \item (Mais difícil) Mostre que $E(X) = 1/\lambda$.
\end{enumerate}
\item Seja $f(x) = Cx^{-3}$ para $x \geq 10$ e zero para $x < 10$.
\begin{enumerate}[label=(\alph*)]
  \item Encontre $C$.
  \item Calcule a função de distribuição acumulada.
  \item Calcule $P(X > 20)$.
\end{enumerate}

\end{enumerate}

\bigskip
\needspace{12\baselineskip}
\section*{Aprofundamento: modelos probabilísticos e dados reais}
\noindent\textit{Exercício ligado à seção de aplicação macroeconômica do Capítulo 4 do livro.}

\begin{enumerate}[label=\textbf{A\arabic*.}, leftmargin=*]
\item \textbf{(Um modelo contínuo para a inflação mensal.)} De janeiro de 2017 a agosto de 2026 (116 meses), o IPCA mensal variou de $-0{,}68\%$ a
$1{,}62\%$, com média $0{,}41\%$; houve 13 meses com inflação negativa e 9 meses acima de 1\% (IBGE, via BCB-SGS). Um analista modela a inflação mensal $X$ (em \%) pela
densidade triangular
$$
f(x) = \begin{cases}
\dfrac{2(x + 0{,}7)}{2{,}53}, & -0{,}7 \leq x \leq 0{,}4, \\[2mm]
\dfrac{2(1{,}6 - x)}{2{,}76}, & 0{,}4 < x \leq 1{,}6, \\[2mm]
0, & \text{caso contrário}.
\end{cases}
$$
\begin{enumerate}[label=(\alph*)]
  \item Verifique que $f$ é uma densidade (use a área de triângulos, sem integrar) e esboce o gráfico.
  \item Obtenha a função de distribuição acumulada $F(x)$ e calcule $P(X < 0)$, $P(X > 1)$ e $F(0{,}4)$.
  \item Calcule $E(X)$. Compare as três probabilidades/médias do modelo com as frequências e a média observadas.
  \item Em que parte da distribuição o modelo erra mais? Que consequência isso teria para um banco que usa o modelo para avaliar o risco de inflação alta?
\end{enumerate}

\end{enumerate}

\end{document}
